
\section{Koszul duality and universal chiral defects}\label{sec:koszul-defects}
\index{quantum Poincare duality@quantum PD|textbf}
\index{modular operad!quantum}
\index{Theorem C!quantum Poincare duality, three tiers}

Poincar\'e duality for manifolds identifies the chain theory with the
linear dual of the cochain theory after an orientation shift.  Quantum
Poincar\'e duality for a chiral algebra is the same assertion after
the ordinary linear dual has been replaced by Verdier duality on the
Ran/Fulton--MacPherson compactification.  The theorem has five
non-interchangeable outputs:
\[
\cA,\qquad
\barBch_X(\cA),\qquad
\cA^{\mathrm i}:=H^\star\barBch_X(\cA),\qquad
\cA^!:=(\cA^{\mathrm i})^\vee,\qquad
Z^{\mathrm{der}}_{\mathrm{ch}}(\cA)
 = C^\bullet_{\mathrm{ch}}(\cA,\cA).
\]
The bar coalgebra classifies twisting morphisms.  The cobar of the
bar coalgebra reconstructs the input:
\[
R_X(\cA):=\Omegach_X\barBch_X(\cA)
\xrightarrow{\sim}\cA
\]
on the Koszul locus and on the weight-completed Positselski surface of
Theorem~B.  The Koszul-dual algebra is produced by the different
functor
\[
K_X(\cA):=\Omegach_X\mathbb{D}_{\Ran}\barBch_X(\cA),
\]
the Verdier-dualized functor of
Theorem~\ref{thm:koszul-reflection}.  Thus
$\Omega\barB(\cA)=\cA$ is reconstruction, not Koszul duality;
$\mathbb{D}_{\Ran}\barB(\cA)$ is the Verdier-dual defect object; and
the cochain-level closed-sector object is the chiral derived centre, not the
bar complex.

The scalar shadow of this Verdier pairing is Theorem~C.  For the
standard witnesses,
\[
\kappa(\cH_k)+\kappa(\cH_k^!)=0,\qquad
\kappa(\widehat{\fg}_k)+\kappa(\widehat{\fg}_k^!)=0,\qquad
\kappa(\mathrm{Vir}_c)+\kappa(\mathrm{Vir}_{26-c})=13.
\]
The landscape census fixes the normalizations
\[
\kappa(V_k(\fg))=\frac{\dim(\fg)(k+h^\vee)}{2h^\vee},\qquad
\kappa(\mathrm{Vir}_c)=\frac c2,\qquad
\kappa(\mathcal W_N)=c(H_N-1).
\]
For $\mathcal W$-families the scalar sum is the
family-dependent Verdier conductor $K^\kappa=\varrho_\cA K$; the
zero-sum law is confined to the Heisenberg, affine, and free-field
lanes listed above.  Beyond the scalar trace, the full quantum correction
is the MC element
\[
\Theta_\cA\in
\MC\!\left(\Defcyc(\cA)\widehat{\otimes}\Gmod\right)
\]
in the completed modular convolution algebra.  Its scalar projection
is $\kappa(\cA)\lambda_g$ at genus~$1$ for all families and at all
genera on the uniform-weight lane; multi-weight algebras at
$g\geq2$ carry the cross-channel correction
$\delta F_g^{\mathrm{cross}}(\cA)$ of
Theorem~\ref{thm:multi-weight-genus-expansion}.  A holographic
comparison may identify the chiral derived centre with a sector of a
$3$d holomorphic-topological target once a separate BV/BRST
comparison has been constructed.  The bar-cobar and Verdier theorem
do not use that comparison as input.

\subsection{Genus-graded Verdier pairing and holographic comparison}

\begin{remark}[Costello--Li comparison across genera]\label{rem:costello-li}
For $N$ D-branes, the genus-graded boundary operator algebra at
$N\to\infty$ and the twisted supergravity BRST algebra are expected
to match after the BV/BRST complex of the physical theory has been
identified with the completed bar-side deformation complex
~\cite{CostelloGaiotto2020,FernandezCostelloP24}.  This is a
holographic comparison theorem once that BV/BRST dictionary is
supplied, not a replacement for the algebraic Verdier construction
of~$\cA^!$.
\end{remark}

\begin{remark}[Genus-graded Verdier pairing]
The universal MC class $\Theta_\cA$
(Theorem~\ref{thm:explicit-theta}) is a point of the completed modular
convolution algebra.  Verdier duality sends the completed bar sector
of~$\cA$ to the completed dual sector of~$\cA^!$ and exchanges the two
Lagrangian summands in the quantum complementarity package.  It does
not assert $\Theta_\cA+\Theta_{\cA^!}=0$.  Its scalar shadow is the
family-dependent Verdier sum of Theorem~C: zero for Heisenberg,
affine Kac--Moody, and free-field pairs, thirteen for the Virasoro
same-family pair $\mathrm{Vir}_c,\mathrm{Vir}_{26-c}$, and
$\varrho_\cA K$ on the broader $\mathcal W$-side.
\end{remark}

\subsection{Universal chiral defects and bar-cobar duality}

\begin{definition}[Universal chiral defect]\label{def:universal-defect}
Let $\mathcal{A}$ be an augmented complete chiral algebra on a smooth
projective curve~$X$ for which the completed bar coalgebra
$\widehat{\barB}^{\mathrm{ch}}_X(\mathcal A)$ is holonomic and
Verdier dualizable on $\Ran(X)$.  The \emph{universal chiral defect
object} is
\[
\mathcal{D}_X(\mathcal A):=
\mathbb{D}_{\Ran}\widehat{\barB}^{\mathrm{ch}}_X(\mathcal A).
\]
It is the Verdier dual of the bar coalgebra.  When it lies in the
domain of the completed chiral cobar functor, its cobar is the
homotopy Koszul-dual algebra
\[
\mathcal A^!_\infty
:=\Omegach_X\mathcal{D}_X(\mathcal A).
\]
On the strict finite-type Koszul surface this reduces to
$\mathcal A^!\simeq (H^\star\barB^{\mathrm{ch}}_X(\mathcal A))^\vee$.
\end{definition}

\begin{theorem}[Verdier defect and Koszul dual {\cite{LV12, Positselski11}}; \ClaimStatusProvedElsewhere]\label{thm:defect-koszul}
Let $\mathcal A$ satisfy the hypotheses of
Definition~\ref{def:universal-defect}, and assume either that
$\mathcal A$ belongs to the strict finite-type Koszul locus or that
the separated completion, finite-dimensional graded quotients,
Mittag--Leffler, proper-support, and continuous-duality hypotheses of
Theorem~\ref{thm:completion-koszul}, together with the convergence
control of Theorem~\ref{thm:bar-convergence}, hold.  Then:
\begin{enumerate}[label=\textup{(\roman*)}]
\item the reconstruction counit is
\[
\Omegach_X\widehat{\barB}^{\mathrm{ch}}_X(\mathcal A)
\xrightarrow{\sim}\mathcal A
\]
in the corresponding ordinary or coderived ambient;
\item the Verdier-dual defect object
$\mathcal{D}_X(\mathcal A)=
\mathbb{D}_{\Ran}\widehat{\barB}^{\mathrm{ch}}_X(\mathcal A)$
is the factorization object whose completed cobar gives the homotopy
Koszul dual $\mathcal A^!_\infty$;
\item on the strict finite-type Koszul surface,
\[
H^\star(\mathcal A^!_\infty)
\simeq
\mathcal A^!
\simeq
\bigl(H^\star\barB^{\mathrm{ch}}_X(\mathcal A)\bigr)^\vee
\simeq
\RHom_{\mathcal A\text{-mod}}(\mathbb C,\mathbb C),
\]
with multiplication induced by Yoneda composition.
\end{enumerate}
Thus $\Omegach_X\barBch_X(\mathcal A)$ recovers~$\mathcal A$, while
$\Omegach_X\mathbb{D}_{\Ran}\barBch_X(\mathcal A)$ produces the
Koszul-dual algebra.  The chiral derived centre
$Z^{\mathrm{der}}_{\mathrm{ch}}(\mathcal A)$ is a Hochschild object
and is not either of these two bar outputs.
\end{theorem}

\subsection{The M2 brane example: Yangian and the Verdier package}

\begin{example}[M2 branes at \texorpdfstring{$A_{N-1}$}{A_N-1} singularity]\label{ex:M2-brane}
For $K$ M2 branes at an $A_{N-1}$ singularity
\cite{Costello-1705.02500v1}, the boundary algebra in the twisted
large-$K$ limit is modeled by a Yangian-type algebra
$Y_{\hbar_{\mathrm{Y}}}(\mathfrak{gl}_N)$, while the proposed
holomorphic-topological bulk target is a completed algebra of
differential operators of the form
$U_{\hbar_{\mathrm{Y}},c}(\mathrm{Diff}(\mathbb C)\otimes\mathfrak{gl}_N)$.
The algebraic theorem supplies the completed filtered bar-cobar and
Verdier package for the Yangian side.  The identification of that
package with the holomorphic-topological target requires the construction of the
higher-spin target and a BV/BRST comparison; it is not a theorem of
ordinary bar-cobar inversion.  The parameter $\hbar_{\mathrm{Y}}$ is
the Yangian deformation parameter, not the genus-counting parameter
used below.
\end{example}

\begin{theorem}[Curved Koszul duality (algebraic form) {\cite{Positselski11, GLZ22}}; \ClaimStatusProvedElsewhere]\label{thm:curved-koszul}
Let $\mathcal{A}$ be an augmented complete filtered chiral algebra with
exhaustive, separated filtration, finite-dimensional graded quotients,
quadratic associated graded algebra, central filtration-$\ge2$
correction classes, and the convergence or completion hypotheses of
Theorems~\ref{thm:filtered-to-curved},~\ref{thm:bar-convergence},
and~\ref{thm:completion-koszul}.
Then:
\begin{enumerate}[label=\textup{(\roman*)}]
\item $\mathcal{A}$ admits a curved $A_\infty$ model with curvature
$m_0\in F_2\mathcal A$ and
\[
m_1^2(x)=m_2(m_0,x)-m_2(x,m_0)=[m_0,x];
\]
\item the completed bar construction is the continuous product
$\widehat{\bar{B}}(\mathcal{A})=\prod_n\bar{B}^n(\mathcal A)$ with
the curved bar differential determined by the completed filtration;
\item Verdier duality yields the curved defect object
\[
\mathcal D_{\mathrm{curv}}(\mathcal A)
=
\mathbb{D}_{\mathrm{Ran}}\widehat{\bar{B}}(\mathcal{A}),
\]
and completed cobar gives the homotopy Koszul-dual algebra
$\mathcal{A}_{\mathrm{curv}}^{!\,\infty}
=\Omegach_X\mathcal D_{\mathrm{curv}}(\mathcal A)$;
\item the completed bar-cobar adjunction gives the
coderived/contraderived Positselski equivalence.  Its reconstruction
part recovers $\mathcal A$ from $\widehat{\bar{B}}(\mathcal A)$; its
Verdier part produces the dual after applying
$\mathbb{D}_{\Ran}$ before cobar.
\end{enumerate}
\end{theorem}


\subsection{\texorpdfstring{The AdS$_3$/CFT$_2$ example: twisted supergravity}{The AdS 3/CFT 2 example: twisted supergravity}}

\begin{example}[AdS\texorpdfstring{$_3 \times S^3 \times T^4$}{_3 x S3 x T4} holography]\label{ex:AdS3}
For type IIB on AdS$_3 \times S^3 \times T^4$ \cite{CP2020}, the
large-$N$ boundary single-trace algebra is a
$W_{1+\infty}$-type algebra at $c=6N$, and the proposed twisted
physical comparison target is Kodaira--Spencer gravity.  The intrinsic algebraic
closed-sector slot attached to the boundary is the Hochschild object
$Z^{\mathrm{der}}_{\mathrm{ch}}(W_{1+\infty})
=C^\bullet_{\mathrm{ch}}(W_{1+\infty},W_{1+\infty})$.
The comparison with KS gravity is the holographic dictionary:
\[
Z^{\mathrm{der}}_{\mathrm{ch}}(W_{1+\infty})
\dashrightarrow
\Obs_{\mathrm{BRST}}(\mathrm{KS\ gravity\ on\ AdS}_3).
\]
\end{example}

\begin{conjecture}[Backreaction as a completed MC perturbation; \ClaimStatusConjectured]\label{conj:backreaction}
\index{gravitational backreaction}
In the AdS$_3$/CFT$_2$ comparison, gravitational backreaction is a
formal deformation of the completed MC element:
\[
\Theta_{1/N}
=
\Theta_0+\sum_{r\ge1}N^{-r}\theta_r
\in
\MC\!\left(\gAmod\widehat{\otimes}\mathbb C[[N^{-1}]]\right).
\]
Writing $D_{\Theta_0}=D+[\Theta_0,-]$, the coefficients satisfy
\[
D_{\Theta_0}\theta_r
+\frac12\sum_{\substack{i+j=r\\i,j\ge1}}[\theta_i,\theta_j]=0.
\]
Generator shifts are $L_\infty$ gauge transformations of this MC
element, and differential shifts are the corresponding twisted
differentials.  A scalar free-energy statement extracted from
$\Theta_{1/N}$ must include the decomposition
\[
F_g(\mathcal A)
=
\kappa(\mathcal A)\lambda_g^{\mathrm{FP}}
+\delta F_g^{\mathrm{cross}}(\mathcal A),
\]
where the cross-channel term vanishes at $g=1$ and on the
uniform-weight lane, and need not vanish for multi-weight
$\mathcal W$-type algebras at $g\ge2$.
\end{conjecture}

A precise physical formulation requires Witten-diagram representatives
for the classes $\theta_r$ and a comparison between the resulting
BRST complex and the completed bar-side complex.  This is the
remaining content of Conjecture~\ref{conj:master-bv-brst}; the
algebraic bar-cobar and Verdier statements do not depend on it.

\begin{theorem}[Finite-type Ext model for the defect algebra {\cite{LV12}}; \ClaimStatusProvedElsewhere]\label{thm:universal-defect-construction}
Assume $\mathcal A$ is augmented, compactly generated, finite type, and
Koszul, and that the augmentation module $\mathbb C$ is compact in the
derived category of $\mathcal A$-modules.  Then the cohomology algebra
of the Verdier defect is computed by the Ext algebra
\[
H^\star\mathcal D_X(\mathcal A)
\simeq
\Ext^\star_{\mathcal A}(\mathbb C,\mathbb C),
\]
with multiplication given by Yoneda product.  A morphism
$\mathcal A\to\mathcal B$ induces the contravariant map
$\Ext^\star_{\mathcal B}(\mathbb C,\mathbb C)
\to
\Ext^\star_{\mathcal A}(\mathbb C,\mathbb C)$ whenever the
augmentation modules and completions are compatible.  The involutive
statement $\mathcal A^{!!}\simeq\mathcal A$ is the finite-type Koszul
involutivity theorem, not a statement about arbitrary defects.
\end{theorem}

\subsection{Examples and computations}

\subsubsection{Example: free fermion and its Koszul dual}

\begin{example}[\texorpdfstring{$bc$}{bc} ghost system \texorpdfstring{$\leftrightarrow$}{iff} \texorpdfstring{$\beta\gamma$}{beta-gamma} system]
The $bc$--$\beta\gamma$ Koszul duality is treated in
\S\ref{sec:fermion-boson-koszul}
(Theorem~\ref{thm:betagamma-bc-koszul}); we record only the
result needed here. The $bc$ ghost system (two fermionic generators
$b, c$ with OPE $b(z)c(w) \sim (z-w)^{-1}$) is Koszul dual to the
bosonic $\beta\gamma$ system:
\[
 bc \text{ ghosts}
 \quad \xleftrightarrow{\text{Koszul}} \quad
 \beta\gamma \text{ system.}
\]
The duality exchanges the exterior (fermionic) relation algebra with
the symmetric (bosonic) one, generalising the classical
$\Lambda(V) \leftrightarrow \mathrm{Sym}(V^*)$ Koszul duality.
\end{example}

\subsubsection{Example: Heisenberg and $\mathcal{W}$-algebras}

\begin{example}[Affine Kac--Moody: Koszul duality vs Feigin--Frenkel involution]
The Feigin--Frenkel involution
$k \mapsto -k - 2h^\vee$ and the chiral Koszul duality
$\widehat{\mathfrak{g}}_k \mapsto (\widehat{\mathfrak{g}}_k)^!$ are
\emph{distinct} operations. The Feigin--Frenkel involution replaces
the level within the \emph{same} family of affine algebras; chiral
Koszul duality produces the Verdier-linear dual of the bar cohomology,
\[
(\widehat{\mathfrak{g}}_k)^!
=
\bigl(H^\star\widehat{\bar{B}}^{\mathrm{ch}}
(\widehat{\mathfrak{g}}_k)\bigr)^\vee
\]
on the finite or completed Koszul surface.  Both the Feigin--Frenkel
partner $\widehat{\mathfrak g}_{-k-2h^\vee}$ and the Koszul dual have
modular characteristic
$-\kappa(\widehat{\mathfrak g}_k)$, because
\[
\kappa(V_k(\fg))
=
\frac{\dim(\fg)(k+h^\vee)}{2h^\vee}.
\]
This scalar agreement does not identify their OPEs, centres, or
Verdier-dual bar models.

For the Heisenberg algebra $\cH_k$, the Koszul dual is
$\cH_k^! \simeq \mathrm{Sym}^{\mathrm{ch}}(V^*)$ with curvature
$-k \cdot \omega_X$ (Part~\ref{part:characteristic-datum}). This is \emph{not} the same as
$\cH_{-k}$: the dual $\cH_k^!$ is generated by $V^*$ with its own OPE
structure, while $\cH_{-k}$ is generated by $V$ with level $-k$. The bar complex
$\bar{B}(\cH_k)$ is the Chevalley--Eilenberg complex
$\mathrm{CE}(\hat{\mathfrak{h}}_k)$; this is the bar construction,
not the Koszul dual itself.

The curvature element $m_0 = \kappa \cdot \omega_X$ satisfies
$m_1^2 = [m_0, -]$ on the curved $A_\infty$ structure.  The
ordinary square-zero differential is obtained only after passing to
the flat representative or to the coderived completed model
(cf.\ Part~\ref{part:characteristic-datum}).
\end{example}

\subsubsection{Yangian from M2 branes}

\begin{calculation}[Yangian structure constants {\cite{Drinfeld85}}; \ClaimStatusProvedElsewhere]
For M2 branes, the Yangian generators $\{E_{ij}^{(r)}\}$ satisfy:

In the Drinfeld ``new'' presentation, the Yangian
$Y_{\hbar_{\mathrm{Y}}}(\mathfrak{gl}_N)$ has generators
$\{T_{ij}^{(r)}\}_{1 \leq i,j \leq N,\, r \geq 0}$ with
$T_{ij}^{(0)} = \delta_{ij}$ and the \emph{RTT relation}:
\[
[T_{ij}^{(r+1)}, T_{kl}^{(s)}]
- [T_{ij}^{(r)}, T_{kl}^{(s+1)}]
=
\hbar_{\mathrm{Y}}\left(T_{kj}^{(r)} T_{il}^{(s)}
- T_{kj}^{(s)} T_{il}^{(r)}\right).
\]

Setting $E_{ij}^{(r)} = T_{ij}^{(r+1)}$ for the shifted generators, the first few relations are:
\begin{align}
[E_{ij}^{(0)}, E_{jk}^{(0)}] &= E_{ik}^{(0)} \quad \text{($\mathfrak{gl}_N$ at level 0)}\\
[E_{ij}^{(0)}, E_{jk}^{(1)}] &= E_{ik}^{(1)} \quad \text{($\mathfrak{gl}_N$-equivariance)}\\
[E_{ij}^{(1)}, E_{jk}^{(1)}] - [E_{ij}^{(0)}, E_{jk}^{(2)}] &= \hbar_{\mathrm{Y}}\left(E_{jj}^{(0)} E_{ik}^{(1)} - E_{jj}^{(1)} E_{ik}^{(0)}\right) \quad \text{(quantum correction)}
\end{align}
\end{calculation}

\subsubsection{\texorpdfstring{Bar complex computation for $\mathcal{W}_3$ algebra}{Bar complex computation for W 3 algebra}}

\begin{example}[\texorpdfstring{$\mathcal{W}_3$}{W_3} bar complex]\label{ex:w3-bar}
For $\mathcal{W}_3$ (the $\mathfrak{sl}_3$ principal W-algebra):

\emph{Generators.} $T$ (spin 2), $W$ (spin 3).

\emph{Bar complex.}
The bar complex is bigraded by bar degree $n$ and conformal weight $h$. At the level of generators (ignoring descendants):
\begin{align}
\dim \bar{B}^0 &= 1 \,\, (\text{vacuum}) \\
\dim \bar{B}^1 &= 2 \,\, (\text{generators: } T, W)
\end{align}
Higher bar degrees involve the $\mathcal{W}_3$ OPE relations. For a \emph{single-generator} operad, the dimensions would follow the Catalan sequence $1, 1, 2, 5, 14, \ldots$; for a two-generator algebra like $\mathcal{W}_3$, the dimensions of $\bar{B}^n_{\leq h}$ (truncated to conformal weight $\leq h$) depend on the specific $\mathcal{W}_3$ OPE relations and require case-by-case computation.
\end{example}

\subsubsection{Critical level phenomena}

\begin{definition}[Critical level]\label{def:critical}
The critical level is $k = -h^\vee$ where $h^\vee$ is the dual Coxeter number. At this level:
\begin{itemize}
\item The Sugawara construction fails (denominator vanishes)
\item The center becomes large (Feigin--Frenkel center)
\item The center identifies with ${}^L\mathfrak{g}$-opers (geometric Langlands)
\end{itemize}
\end{definition}

\begin{theorem}[Feigin--Frenkel center {\cite{Feigin-Frenkel,BD04}}; \ClaimStatusProvedElsewhere]\label{thm:ff-center}
At critical level, the center of $\widehat{\mathfrak{g}}_{-h^\vee}$ is:
\[
Z(\widehat{\mathfrak{g}}_{-h^\vee}) \cong \mathrm{Fun}(\mathrm{Op}_{{}^L\mathfrak{g}}(D^\times))
\]
functions on the space of ${}^L\mathfrak{g}$-opers on the punctured formal disk~$D^\times$. The global version on a curve~$X$ follows from Beilinson--Drinfeld's chiral algebra formalism~\cite{BD04}.
\end{theorem}

\begin{remark}[Opers and connections]
An oper is a special kind of connection:
\[
\nabla = \partial + p_{-1} + \text{regular terms}
\]
where $p_{-1}$ is a principal nilpotent element. Opers are related to solutions of the KZ equations through the theory of conformal blocks.
\end{remark}

\subsubsection{\texorpdfstring{Chiral coalgebra structure for $\beta\gamma$}{Chiral coalgebra structure for beta-gamma}}

\begin{proposition}[\texorpdfstring{$\beta\gamma$}{beta-gamma} bar complex coalgebra; \ClaimStatusProvedHere]\label{prop:bg-bar-coalg}
The bar complex $\bar{B}^{\text{ch}}(\beta\gamma)$ has chiral coalgebra structure:
\begin{enumerate}
\item \emph{Comultiplication}: Elements decompose as:
\[
\Delta(\beta_{i_1} \cdots \beta_{i_p} \gamma_{j_1} \cdots \gamma_{j_q} \partial^k) =
\sum_{\substack{I_\beta \sqcup I'_\beta = \{i_1,\ldots,i_p\} \\ I_\gamma \sqcup I'_\gamma = \{j_1,\ldots,j_q\}}}
\beta_{I_\beta}\gamma_{I_\gamma}\partial^{k_1} \otimes \beta_{I'_\beta}\gamma_{I'_\gamma}\partial^{k_2}
\]
respecting normal ordering: $\beta$'s to the left of $\gamma$'s.

\item \emph{Growth:} The bar complex is bigraded by bar degree $n$ and conformal weight $h$. At each conformal weight, $\bar{B}^n_h$ is finite-dimensional with generating function
\[\sum_{n,h} \dim(\bar{B}^n_h)\, q^h t^n = \prod_{m \geq 0} \frac{1}{(1-q^{m+\lambda}t)(1-q^{m+1-\lambda}t)}\]
reflecting the two families of generators $\beta_{-m-\lambda}$ and $\gamma_{-m-1+\lambda}$ in the augmentation ideal. Without a conformal weight bound, $\dim(\bar{B}^n) = \infty$ for each $n \geq 1$.

\item \emph{Coassociativity:} Follows from the factorization property of the FM operad. The boundary divisors $D_S \subset \overline{C}_{n}(X)$ parametrizing collision of a subset $S$ of points induce maps
\[
\overline{C}_{n}(X) \supset D_S \;\cong\; \overline{C}_{|S|}(X) \times \overline{C}_{n-|S|+1}(X)
\]
and these compose associatively by the operad axioms.
\end{enumerate}
\end{proposition}

\begin{proof}
Part~(3) is a special case of Theorem~\ref{thm:bar-cobar-isomorphism-main}: the bar complex of \emph{any} augmented chiral algebra carries a factorization coalgebra structure, and coassociativity is the operad axiom for the boundary strata of~$\overline{C}_n(X)$. Part~(1) is the explicit formula for the \emph{factorization coproduct} on $\bar{B}^{\mathrm{ch}}(\beta\gamma)$: the chiral bracket of the $\beta\gamma$ system is bilinear in modes, so when two groups of points separate, a normally-ordered monomial $\beta_{i_1}\cdots\beta_{i_p}\gamma_{j_1}\cdots\gamma_{j_q}$ decomposes by distributing its mode indices between the two groups. This is the Wick-type decomposition inherited from the free-field OPE; the normal-ordering convention is preserved because the factorization coproduct respects the weight grading. Part~(2) is a character computation: the augmentation ideal $\overline{\beta\gamma}$ has generators of conformal weights $m + \lambda$ ($\beta$ modes) and $m + 1 - \lambda$ ($\gamma$ modes) for $m \geq 0$, so the bar complex at bar degree~$n$ and weight~$h$ counts $n$-fold tensor products of these generators with total weight~$h$, giving the stated plethystic product. The infinite-dimensionality of~$\bar{B}^n$ at unfixed weight is standard for any chiral algebra with infinitely many modes; the positive-energy axiom ensures finite-dimensionality at each weight.
\end{proof}

\begin{remark}[Physical interpretation]
The coalgebra structure arises from correlation functions on punctured curves.

\emph{Propagator expansion.} The $\beta\gamma$ propagator:
\[
\langle \beta(z)\gamma(w) \rangle = \frac{1}{z-w}
\]
defines a distribution on $C_2(X) = X \times X \setminus \Delta$.

\emph{Tree factorisation.} Higher correlations factor through tree graphs:
\[
\langle \beta(z_1)\gamma(z_2)\beta(z_3)\gamma(z_4) \rangle =
\sum_{\text{pairings}} \prod_{\text{edges}} \frac{1}{z_i - z_j}
\]

\emph{Fulton--MacPherson compactification.} The Fulton--MacPherson compactification $\overline{C}_n(X)$
regularizes these distributions, with the coalgebra structure encoding how correlators
factorize when points collide.
\end{remark}

\subsection{The prism principle in action}

\begin{example}[Structure coefficients via residues]
For generators $\phi_i$ with OPE $\phi_i(z) \phi_j(w) = \sum_k C_{ij}^k \phi_k(w)/(z-w)^{h_i + h_j - h_k} + \cdots$, the bar complex extracts $\mathrm{Res}_{D_{ij}}[\phi_i \otimes \phi_j \otimes \eta_{ij}] = \sum_k C_{ij}^k \phi_k$.
\end{example}

\subsection{The Ayala--Francis perspective}

\begin{theorem}[Factorization homology and the bar complex {\cite{Francis2013,HA}}; \ClaimStatusProvedElsewhere]\label{thm:fact-homology-quantum}
For a chiral algebra $\mathcal{A}$ on $X$, the Ayala--Francis factorization homology and the geometric bar complex are related by Koszul duality:
\[\int_X \mathcal{A} \simeq \mathcal{A} \otimes^{\mathbb{L}}_{\mathrm{Disk}(X)} \mathrm{pt} \qquad \text{(chains)}, \qquad \bar{B}^{\mathrm{ch}}(\mathcal{A}) \simeq \mathrm{RHom}_{\mathcal{A}\text{-mod}}(k, k) \qquad \text{(cochains)}\]
These are \emph{dual} constructions: factorization homology computes the derived tensor product (the \emph{homological} side), while the derived Hom (the \emph{cohomological} side) is computed via the bar complex. They are related by $\mathbb{E}_2$-Koszul duality as developed by Francis \cite{Francis2013}, building on Lurie's operadic bar-cobar framework~\cite[§5.2]{HA}.
\end{theorem}

\begin{proof}
By Francis \cite{Francis2013}, both constructions are derived functors of the forgetful functor from $\mathcal{A}$-modules to chain complexes, evaluated on the trivial module $k$. The factorization homology $\int_X \mathcal{A} \simeq \mathcal{A} \otimes^{\mathbb{L}}_{\mathrm{Disk}(X)} \mathrm{pt}$ is the derived tensor product (left adjoint, homological), while the bar complex $\bar{B}^{\mathrm{ch}}(\mathcal{A}) \simeq \mathrm{RHom}_{\mathcal{A}\text{-mod}}(k, k)$ is the derived Hom (right adjoint, cohomological). These are related by $\mathbb{E}_2$-Koszul duality: the $\mathbb{E}_2$-operad is Koszul self-dual up to shift \cite[Theorem~5.2]{Francis2013}, which induces the duality between the two functors.
\end{proof}

\subsection{The prism principle: operadic identification}%
\label{subsec:prism-operadic}

The geometric bar-cobar duality of chiral algebras is an instance of
\emph{operadic Koszul duality} in the sense of Ginzburg--Kapranov
\cite{GK94} and Getzler \cite{Get95}.

\subsubsection{The chiral operad}

\begin{definition}[Chiral operad]\label{def:chiral-operad}
The \emph{chiral operad} $\mathcal{P}_{\mathrm{ch}}$ is the topological
operad with:
\[
\mathcal{P}_{\mathrm{ch}}(n)
= \Omega^*_{\log}\!\bigl(\overline{C}_{n+1}(X)\bigr),
\qquad n \geq 1,
\]
the complex of logarithmic differential forms on the Fulton--MacPherson
compactification of $n+1$ points on the curve $X$. The operad
composition maps:
\[
\circ_i \colon \mathcal{P}_{\mathrm{ch}}(m) \otimes
\mathcal{P}_{\mathrm{ch}}(n) \to \mathcal{P}_{\mathrm{ch}}(m+n-1)
\]
are induced by the gluing maps
$\overline{C}_{m+1}(X) \times \overline{C}_{n+1}(X) \to
\overline{C}_{m+n}(X)$ that identify the $i$-th point of the first
configuration with the distinguished point of the second.

The symmetric group $\mathfrak{S}_n$ acts by permuting the
non-distinguished points, and the unit is the identity form in
$\Omega^0(\overline{C}_2(X)) = \mathbb{C}$.
\end{definition}

\begin{proposition}[Genus-zero identification; \ClaimStatusProvedHere]%
\label{prop:chiral-operad-genus0}
At genus $0$ (i.e., $X = \mathbb{P}^1$), the chiral operad is
quasi-isomorphic to the hypercommutative operad:
\[
\mathcal{P}_{\mathrm{ch}}|_{g=0} \;\simeq\; \mathrm{HyCom}.
\]
Explicitly, $\overline{C}_{n+1}(\mathbb{P}^1) \cong
\overline{\mathcal{M}}_{0,n+2}$ (the Deligne--Mumford compactification
of the moduli space of rational curves with $n+2$ marked points), and
the cohomology operad
$\{H^*(\overline{\mathcal{M}}_{0,n+2})\}_{n \geq 1}$ is the
hypercommutative operad $\mathrm{HyCom}$ of Getzler \cite{Get95}.
\end{proposition}

\begin{proof}
The identification
$\overline{C}_{n+1}(\mathbb{P}^1) \cong \overline{\mathcal{M}}_{0,n+2}$
is classical: the Fulton--MacPherson compactification of $n+1$ points on
$\mathbb{P}^1$ with one point fixed at $\infty$ is isomorphic to the
Deligne--Mumford moduli space of stable rational curves with $n+2$
punctures \cite{Knudsen83}. Under this identification, the operad composition
maps of $\mathcal{P}_{\mathrm{ch}}$ correspond to the clutching maps
\[
\overline{\mathcal{M}}_{0,m+1} \times
\overline{\mathcal{M}}_{0,n+1} \to
\overline{\mathcal{M}}_{0,m+n}
\]
that glue two rational curves at marked points. The cohomology operad
of $\{\overline{\mathcal{M}}_{0,n+2}\}$ is the hypercommutative operad
by Getzler's theorem \cite[Theorem~1.5]{Get95}: the fundamental classes
of the boundary strata satisfy exactly the relations defining
$\mathrm{HyCom}$.

The quasi-isomorphism follows because
$\overline{\mathcal{M}}_{0,n+2}$ is a smooth projective variety
with pure Hodge structure, hence formal as a dga by
Deligne--Griffiths--Morgan--Sullivan~\cite{DGMS75}. Thus
$\Omega^*(\overline{\mathcal{M}}_{0,n+2})$ is quasi-isomorphic
to $H^*(\overline{\mathcal{M}}_{0,n+2})$ as an operad.
\end{proof}

\subsubsection{Operadic Koszul duality and the prism}

\begin{theorem}[Prism principle: operadic identification;
\ClaimStatusProvedHere]\label{thm:prism-operadic}
\textup{[Regime: quadratic at genus~$0$; curved-central at genus~$g \geq 1$
\textup{(}Convention~\textup{\ref{conv:regime-tags})}.]}

The geometric bar-cobar duality of chiral algebras
(Definition~\ref{def:geometric-bar},
Definition~\ref{def:geom-cobar-intrinsic}) is an instance of operadic
Koszul duality. Precisely:
\begin{enumerate}[label=\textup{(\roman*)}]
\item \emph{Genus~$0$:} The chiral operad
 $\mathcal{P}_{\mathrm{ch}}|_{g=0} \simeq \mathrm{HyCom}$ is Koszul,
 with Koszul dual cooperad:
 \[
 \mathcal{P}_{\mathrm{ch}}^{\text{\textexclamdown}}|_{g=0}
 \;\simeq\; \mathrm{Grav}^c,
 \]
 the gravity cooperad. Here $\mathrm{Grav}(n) =
 H_{n-2}(\mathcal{M}_{0,n+2})$ is the gravity operad
 of Getzler \cite{Get95} (note: the \emph{open} moduli space
 $\mathcal{M}_{0,n+2}$, not the compactification $\overline{\mathcal{M}}_{0,n+2}$), and $\mathrm{Grav}^c$ is its cooperad
 incarnation.

\item \emph{Bar as operadic bar:} For a chiral algebra $\mathcal{A}$
 that is an algebra over $\mathcal{P}_{\mathrm{ch}}$, the geometric
 bar complex is the operadic bar construction:
 \[
 \bar{B}^{\mathrm{ch}}(\mathcal{A})
 \;\simeq\;
 B_{\mathcal{P}_{\mathrm{ch}}}(\mathcal{A})
 := \mathcal{P}_{\mathrm{ch}}^{\text{\textexclamdown}}
 \circ_{\mathcal{P}_{\mathrm{ch}}} \mathcal{A}.
 \]
 The bar differential $d_{\mathrm{bar}} = d_{\mathrm{int}} +
 d_{\mathrm{res}} + d_{\mathrm{config}}$ corresponds to the
 standard operadic bar differential: $d_{\mathrm{res}}$ is the
 twisting cochain component and $d_{\mathrm{config}}$ is the
 cooperad cocomposition component.

\item \emph{Cobar as operadic cobar:} Dually, the geometric cobar
 complex is the operadic cobar construction:
 \[
 \Omega^{\mathrm{ch}}(\mathcal{C})
 \;\simeq\;
 \Omega_{\mathcal{P}_{\mathrm{ch}}^{\text{\textexclamdown}}}(\mathcal{C})
 := \mathcal{P}_{\mathrm{ch}}
 \circ_{\mathcal{P}_{\mathrm{ch}}^{\text{\textexclamdown}}} \mathcal{C}.
 \]
 Under Verdier duality
 (Definition~\ref{def:geom-cobar-intrinsic}), the gravity cooperad
 $\mathrm{Grav}^c$ controlling the bar maps to the hypercommutative
 operad $\mathrm{HyCom}$ controlling the cobar.

\item \emph{Bar-cobar resolution:} For a Koszul chiral algebra
 $\mathcal{A}$ (Definition~\ref{def:chiral-koszul-pair}), the
 augmentation map
 $\varepsilon\colon
 \Omega^{\mathrm{ch}}(\bar{B}^{\mathrm{ch}}(\mathcal{A}))
 \to \mathcal{A}$
 (Proposition~\ref{prop:cobar-bar-augmentation})
 is a quasi-isomorphism. This is an instance of the general operadic
 bar-cobar resolution for Koszul operads \cite[Theorem~11.3.3]{LV12}.
\end{enumerate}
\end{theorem}

\begin{proof}
\emph{Part (i): HyCom--Grav Koszul duality.}
This is Getzler's theorem \cite[Theorem~2.3]{Get95}. The
hypercommutative operad $\mathrm{HyCom}$ is defined by the
presentation: generators are the fundamental classes
$[\overline{\mathcal{M}}_{0,n+2}] \in H^{2(n-1)}(
\overline{\mathcal{M}}_{0,n+2})$ with the WDVV (associativity)
relations arising from boundary divisors. The Koszul dual cooperad
has components $\mathrm{Grav}^c(n) = H_{n-2}(
\mathcal{M}_{0,n+2})^{\vee}$, which identifies with the gravity
cooperad (the linear dual of Getzler's gravity operad
$\mathrm{Grav}(n) = H_{n-2}(\mathcal{M}_{0,n+2})$, defined using
the \emph{open} moduli space). The Koszulity of $\mathrm{HyCom}$ is
established in \cite[Theorem~2.3]{Get95} via the purity of the Hodge
structure on $H^*(\overline{\mathcal{M}}_{0,n})$.

\emph{Part (ii): Identification of the bar complex.}
The geometric bar complex (Definition~\ref{def:geometric-bar}) at
genus~$0$ is:
\[
\bar{B}^{\mathrm{ch}}_n(\mathcal{A}) = \Gamma\!\left(
 \overline{C}_{n+1}(\mathbb{P}^1),\,
 j_*j^*\mathcal{A}^{\boxtimes(n+1)}
 \otimes \Omega^n_{\log}
\right).
\]
The logarithmic forms $\Omega^*_{\log}(\overline{C}_{n+1}(\mathbb{P}^1))$
compute the cohomology of the \emph{open} moduli space
$\mathcal{M}_{0,n+2}$ by the logarithmic comparison theorem (Deligne),
providing the propagators for the bar differential. The bar complex
is controlled by the cohomology operad of
$\overline{\mathcal{M}}_{0,n+2}$, i.e., $\mathrm{HyCom}$.

The standard operadic bar construction
$B_{\mathrm{HyCom}}(\mathcal{A}) = \mathrm{Grav}^c \circ_{\mathrm{HyCom}}
\mathcal{A}$ is a two-sided composition product. Concretely, its $n$-th
component is:
\[
B_{\mathrm{HyCom}}(\mathcal{A})_n
= \bigoplus_{\text{trees } T} \bigotimes_{v \in V(T)}
 \mathrm{Grav}^c(|v|) \otimes_{\mathfrak{S}_{|v|}} \mathcal{A}^{\otimes |v|}
\]
summed over planar trees $T$ with $n$ leaves. The three components of
the bar differential correspond to:
\begin{itemize}
\item $d_{\mathrm{res}}$: the twisting cochain
 $\kappa\colon \mathrm{Grav}^c \to \mathrm{HyCom}$
 (Ginzburg--Kapranov \cite[§2.1]{GK94}), which maps
 the gravity cooperad to the hypercommutative operad. Geometrically,
 this is the residue map: extracting the leading OPE coefficient at
 a collision divisor.
\item $d_{\mathrm{config}}$: the cooperad cocomposition
 $\Delta\colon \mathrm{Grav}^c \to \mathrm{Grav}^c \circ
 \mathrm{Grav}^c$, which geometrically corresponds to the de Rham
 differential on configuration forms and the
 boundary stratification of $\overline{\mathcal{M}}_{0,n}$.
\item $d_{\mathrm{int}}$: the internal differential of $\mathcal{A}$.
\end{itemize}

\emph{Part (iii): Cobar via Verdier duality.}
By the intrinsic cobar definition
(Definition~\ref{def:geom-cobar-intrinsic}), $\Omega^{\mathrm{ch}}$
is the Verdier dual of the bar. Since Verdier duality on
$\overline{\mathcal{M}}_{0,n+2}$ exchanges $\mathrm{Grav}^c$ and
$\mathrm{HyCom}$ (this is the geometric content of Koszul duality:
the Koszul dual operad is the Poincaré dual of the Koszul dual
cooperad), the operadic cobar construction
$\Omega_{\mathrm{Grav}^c}(\mathcal{C}) = \mathrm{HyCom}
\circ_{\mathrm{Grav}^c} \mathcal{C}$ agrees with the
geometric cobar $\Omega^{\mathrm{ch}}(\mathcal{C})$.

\emph{Part (iv): Bar-cobar resolution.}
For a Koszul operad $\mathcal{P}$, the canonical map
$\Omega(B_{\mathcal{P}}(\mathcal{A})) \to \mathcal{A}$ is a
quasi-isomorphism for any $\mathcal{P}$-algebra $\mathcal{A}$ (this
is the fundamental theorem of operadic Koszul duality
\cite[Theorem~11.3.3]{LV12}). Since $\mathrm{HyCom}$ is Koszul
by Part~(i), and the geometric bar-cobar agree with the operadic
bar-cobar by Parts~(ii)--(iii), the result follows.
\end{proof}

\begin{definition}[Modular operad]\label{def:modular-operad}
\index{modular operad|textbf}
A \emph{modular operad} (Getzler--Kapranov \cite[\S5]{GeK98})
in a symmetric monoidal category $\mathcal{V}$ consists of:
\begin{enumerate}[label=(\roman*)]
\item A collection of objects
 $\{\mathcal{M}(g,n)\}_{2g-2+n>0}$ in $\mathcal{V}$,
 each equipped with a $\Sigma_n$-action (the stability condition
 $2g-2+n > 0$ excludes $(g,n) = (0,0), (0,1), (0,2)$).
\item \emph{Separating gluing maps}: for each partition
 $\{1, \ldots, n\} = S_1 \sqcup S_2$ and $g = g_1 + g_2$,
 \[
 \xi_{S_1, S_2}\colon \mathcal{M}(g_1, |S_1|+1)
 \otimes \mathcal{M}(g_2, |S_2|+1)
 \;\longrightarrow\; \mathcal{M}(g, n),
 \]
 corresponding to gluing two curves at marked points.
\item \emph{Non-separating gluing maps}: for each pair of
 distinguished marked points,
 \[
 \xi_{\mathrm{ns}}\colon \mathcal{M}(g, n+2)
 \;\longrightarrow\; \mathcal{M}(g+1, n),
 \]
 corresponding to self-gluing (identifying two marked points
 on the same curve).
\end{enumerate}
These maps satisfy:
\begin{itemize}
\item \emph{Equivariance}: compatibility with the $\Sigma_n$-actions;
\item \emph{Associativity}: iterated gluings are independent of
 the order in which they are performed;
\item \emph{Commutativity}: separating gluing is symmetric in
 the two factors.
\end{itemize}
The canonical example is the collection of Deligne--Mumford moduli
spaces $\{\overline{\mathcal{M}}_{g,n}\}_{2g-2+n>0}$ with clutching
maps at the boundary strata.
\end{definition}

\begin{definition}[Feynman transform]\label{def:feynman-transform}
\index{Feynman transform|textbf}
Let $\mathcal{M}od$ be a modular operad and $\mathcal{A}$ an algebra
over $\mathcal{M}od$. The \emph{Feynman transform}
$\mathrm{FT}_{\mathcal{M}od}(\mathcal{A})$
\cite[Definition~4.1]{GeK98} is the chain complex
\[
\mathrm{FT}_{\mathcal{M}od}(\mathcal{A})
\;=\;
\bigoplus_{\substack{g,n \\ 2g-2+n>0}}
\;\bigoplus_{\Gamma \in \mathcal{G}_{g,n}}
\bigotimes_{v \in V(\Gamma)} \mathcal{A}(\mathrm{val}(v))
\;\otimes\;
\bigotimes_{e \in E(\Gamma)} \det(e),
\]
summed over \emph{stable graphs} $\Gamma$ of type $(g,n)$
(connected graphs with genus decoration $g_v$ at each vertex $v$
satisfying $\sum_v g_v + b_1(\Gamma) = g$ and
$\sum_v \mathrm{val}(v) = n + 2|E(\Gamma)|$, where $b_1$ is the
first Betti number and $\mathrm{val}(v)$ counts half-edges at~$v$).
The differential is the signed sum over edge contractions:
\begin{equation}\label{eq:feynman-transform-differential}
d_{\mathrm{FT}}
\;=\;
\sum_{\substack{e \in E(\Gamma) \\ e\;\text{internal}}}
(-1)^{\epsilon(e,\Gamma)}\,
\xi_e^*,
\end{equation}
where $\xi_e^*$ contracts the edge~$e$ of~$\Gamma$ (merging
its two endpoint vertices and applying the modular operad
composition at the contracted marking pair), and
$\epsilon(e, \Gamma) = \sum_{e' < e} 1$ is the position
of~$e$ in the chosen total ordering of edges.
More precisely, each edge~$e$ connecting vertices $v_1, v_2$
at half-edge labels $i_1, i_2$ contributes
\[
\xi_e^*\colon
\cA(\mathrm{val}(v_1)) \otimes \cA(\mathrm{val}(v_2))
\;\xrightarrow{\;\xi_{i_1 i_2}\;}
\cA(\mathrm{val}(v_1) + \mathrm{val}(v_2) - 2),
\]
using the separating gluing map of the modular operad.
For a self-loop at vertex~$v$ connecting half-edges $i_1, i_2$,
the contribution is the non-separating gluing
$\xi_{\mathrm{ns}}\colon \cA(\mathrm{val}(v)) \to
\cA(\mathrm{val}(v) - 2)$.
The identity $d_{\mathrm{FT}}^2 = 0$ follows because each
codimension-$2$ face (edge pair $\{e, e'\}$) appears in
exactly two codimension-$1$ faces with opposite signs
$(-1)^{\epsilon(e,\Gamma)} \cdot (-1)^{\epsilon(e',\Gamma/e)}
+ (-1)^{\epsilon(e',\Gamma)} \cdot
(-1)^{\epsilon(e,\Gamma/e')} = 0$,
the standard boundary-of-boundary cancellation.
\end{definition}

\begin{theorem}[Prism principle: higher-genus extension;
\ClaimStatusProvedHere]\label{thm:prism-higher-genus}
The operadic identification extends to all genera via the
modular operad structure on $\{\overline{\mathcal{M}}_{g,n}\}$
(Definition~\ref{def:modular-operad}).
Precisely:
\begin{enumerate}[label=\textup{(\roman*)}]
\item The collection $\mathcal{M}od = \{\overline{\mathcal{M}}_{g,n}
 \}_{2g-2+n>0}$ forms a modular operad in the sense of
 Getzler--Kapranov \cite{GeK98}, with composition maps given by
 clutching at marked points.
\item The full genus-graded bar complex
 $\bar{B}^{\mathrm{full}}(\mathcal{A}) = \bigoplus_{g \geq 0}
 \bar{B}^{(g)}(\mathcal{A})$ is the \emph{Feynman transform}
 of $\mathcal{A}$ with respect to the modular operad
 $\mathcal{M}od$:
 \[
 \bar{B}^{\mathrm{full}}(\mathcal{A}) \simeq
 \mathrm{FT}_{\mathcal{M}od}(\mathcal{A}).
 \]
\item The nilpotence $d_{\mathrm{bar}}^2 = 0$ at all genera
 (Theorem~\ref{thm:genus-induction-strict}) is a consequence of the
 modular operad boundary relation $\partial^2 = 0$ on the
 stratification of $\overline{\mathcal{M}}_{g,n}$, combined with the
 centrality of the curvature $\mu_0$.
\end{enumerate}
\end{theorem}

\begin{proof}
\emph{Part (i)} is the foundational result of
Getzler--Kapranov \cite[§5]{GeK98}: the boundary stratification of
$\overline{\mathcal{M}}_{g,n}$ gives gluing maps
\[
\xi_{g_1,g_2}\colon \overline{\mathcal{M}}_{g_1,n_1+1} \times
\overline{\mathcal{M}}_{g_2,n_2+1} \to
\overline{\mathcal{M}}_{g_1+g_2,n_1+n_2}
\quad (\text{separating})
\]
and self-gluing maps
$\xi_g\colon \overline{\mathcal{M}}_{g,n+2} \to
\overline{\mathcal{M}}_{g+1,n}$
(non-separating). These satisfy the axioms of a modular operad.

\emph{Part (ii).}
The Feynman transform (Definition~\ref{def:feynman-transform})
applied to the modular operad
$\mathcal{M}od = \{\overline{\mathcal{M}}_{g,n}\}$ yields a sum
controlled by the
boundary strata of $\overline{\mathcal{M}}_{g,n}$, which are
indexed by stable graphs. The genus-$g$ bar complex
$\bar{B}^{(g)}(\mathcal{A})$ integrates
$\mathcal{A}^{\boxtimes n}$ against logarithmic forms on
$\overline{\mathcal{M}}_{g,n}$, which is precisely what
the Feynman transform produces: each stable graph $\Gamma$
contributes an integral over the corresponding boundary stratum,
with the logarithmic propagator $\eta_{ij}$ along each edge.

\emph{Part (iii).}
The Feynman transform differential squares to zero by the
modular operad axioms: the modular operad boundary relation states
that the composition $\partial \circ \partial$ on the
stratification of $\overline{\mathcal{M}}_{g,n}$ vanishes (each
codimension-2 stratum is reached from exactly two codimension-1
strata with opposite orientations). This is exactly the mechanism
proved in Theorem~\ref{thm:genus-induction-strict}: the Arnold
relations at genus~$0$, the centrality of the curvature $\mu_0$
for cross-genus terms, and the modular operad boundary relation
for pure higher-genus terms combine to give
$d_{\mathrm{bar}}^2 = 0$.
\end{proof}

\begin{corollary}[Genus-$0$ reduction to the operadic bar construction;
\ClaimStatusProvedHere]
\label{cor:feynman-genus0-reduction}
\index{Feynman transform!genus-zero reduction}
The genus-$0$ truncation of the Feynman transform recovers the
operadic bar construction:
\[
\mathrm{FT}_{\mathcal{M}od}(\cA)\big|_{g=0}
\;\simeq\;
\barB_{\mathrm{op}}(\cA).
\]
\end{corollary}

\begin{proof}
At genus~$0$, the only stable graphs are trees
($b_1(\Gamma) = 0$, each vertex has $g(v) = 0$, stability
$2g(v) - 2 + \mathrm{val}(v) > 0$ forces $\mathrm{val}(v) \ge 3$).
Trees with $n$ leaves and internal vertices of valence $\ge 3$ are
precisely the objects indexing the operadic bar construction. The
edge determinant $\det(E(\Gamma))$ provides the correct suspension
signs. The Feynman-transform differential restricted to trees
contracts one internal edge, the operadic bar differential. No
loop graphs, no genus labels, no $\hbar_{\mathrm{g}}$-weighting.
\end{proof}

\begin{corollary}[The loop expansion is the genus expansion;
\ClaimStatusProvedHere]
\label{cor:hbar-genus-identification}
\index{loop expansion!genus identification}
\index{genus expansion!loop identification}
Assign to each stable graph $\Gamma$ the weight
$\hbar_{\mathrm{g}}^{g(\Gamma)}$, where
$g(\Gamma) = b_1(\Gamma) + \sum_v g(v)$ is the total genus.
Then the Feynman transform decomposes as
\[
\mathrm{FT}_{\mathcal{M}od}(\cA)
\;=\;
\sum_{g \ge 0} \hbar_{\mathrm{g}}^g\;
\mathrm{FT}_{\mathcal{M}od}(\cA)\big|_{\text{genus } g},
\]
and $g(\Gamma)$ equals the number of independent loops in~$\Gamma$
plus the sum of the genus decorations at its vertices. At
$\hbar_{\mathrm{g}} = 0$ (genus~$0$), only trees survive
(Corollary~\textup{\ref{cor:feynman-genus0-reduction}}). The
$\hbar_{\mathrm{g}}$-expansion is the genus expansion: each power
of~$\hbar_{\mathrm{g}}$ counts one unit of topological complexity
added to the underlying surface, whether by forming a loop in the graph
($b_1 \to b_1 + 1$) or by increasing a vertex genus
($g(v) \to g(v) + 1$).
\end{corollary}

\begin{proof}
The genus grading on stable graphs is additive under edge
contraction and vertex splitting, hence compatible with the
Feynman-transform differential. The
identification $g(\Gamma) = b_1(\Gamma) + \sum_v g(v)$ is the
definition of total genus for a stable graph
(Definition~\ref{def:stable-graph}).
\end{proof}

\begin{remark}[\texorpdfstring{$\hbar$}{hbar} parameters]
The parameter $\hbar_{\mathrm{g}}$ in
Corollary~\ref{cor:hbar-genus-identification} is a formal genus
counter.  It is distinct from the Yangian parameter
$\hbar_{\mathrm{Y}}$ in the M2 example and from the BV loop-counting
parameter $\hbar_{\mathrm{BV}}$ of
Remark~\ref{rem:bvbrst-two-hbar}.  No numerical specialization of
one parameter is used to define another.
\end{remark}

\begin{example}[The Heisenberg genus expansion from graph counting]
\label{ex:heisenberg-genus-from-graphs}
\index{Heisenberg algebra!genus expansion from graphs}
For the Heisenberg algebra $\cH_k$, every transferred operation
$\Phi_v^\cA$ at a vertex of valence $\ge 3$ vanishes (the OPE is
abelian). The only stable graphs that contribute to
$\Theta_{\cH_k}$ are therefore graphs with all vertices bivalent.
Stability $2g(v) - 2 + \mathrm{val}(v) > 0$ with
$\mathrm{val}(v) = 2$ forces $g(v) \ge 1$. The contributing
graphs at genus~$g$ with no external legs are all connected
graphs on bivalent genus-$1$ vertices. The edge-vertex relation
$2|E| = \sum \mathrm{val}(v) = 2|V|$ gives $|E| = |V|$, so
every connected component is a single cycle. Connectedness
forces a single cycle of length~$|V|$. Its genus is $g(\Gamma) = b_1 + \sum g(v) = 1 +
(g-1) = g$. At each edge, the contraction inserts the metric
$k = \langle\alpha,\alpha\rangle$, and at each vertex the
genus-$1$ operation contributes a factor from the self-trace.
The symmetry factor is $|\mathrm{Aut}(\Gamma)| = 2g$ (rotations
and reflections of a $g$-gon). The resulting coefficient is
\[
F_g(\cH_k)
\;=\;
\frac{k^g}{2g}\cdot
\operatorname{tr}\bigl(\underbrace{\Delta \circ \cdots
\circ \Delta}_{g}\bigr)
\;=\;
k\cdot\frac{2^{2g-1}-1}{2^{2g-1}}\cdot\frac{|B_{2g}|}{(2g)!},
\]
where $B_{2g}$ is the Bernoulli number, via the
Faber--Pandharipande $\lambda_g$~formula
\cite{FP03}. The genus expansion
$\sum_{g \ge 1} F_g\,x^{2g} = k\bigl(\hat{A}(ix) -
1\bigr)$ recovers the $\hat{A}$-genus from graph counting.
\end{example}

\begin{definition}[Modular graph coefficient algebra]
\label{def:modular-graph-algebra}
\index{modular graph coefficient algebra|textbf}
\index{Feynman transform!coefficient algebra}
The \emph{modular graph coefficient algebra} is the completed dg Lie
algebra
\begin{equation}\label{eq:gmod-def}
\Gmod
\;:=\;
\mathrm{FT}_{\mathcal{M}od}\!\bigl(
 \{R\Gamma(\overline{\mathcal{M}}_{g,n},\, \mathbb{Q})\}_{2g-2+n>0}
\bigr),
\end{equation}
the Feynman transform of the modular operad
$\mathcal{M}od = \{\overline{\mathcal{M}}_{g,n}\}$
(Definition~\ref{def:modular-operad}) applied to the rational
cohomology of the moduli spaces themselves. The Lie bracket is
induced by edge-gluing of stable graphs; the differential is induced
by boundary stratification. Each stable graph~$\Gamma$ contributes
the integral over the corresponding modular stratum, weighted by the
automorphism factor~$|\mathrm{Aut}(\Gamma)|^{-1}$.

\smallskip\noindent
\emph{Role.}
The universal Maurer--Cartan class
(Theorem~\ref{thm:universal-MC}) lives in
$\operatorname{MC}(\Defcyc(\cA) \;\widehat{\otimes}\; \Gmod)$:
the cyclic deformation complex of~$\cA$ tensored with the modular
graph coefficient algebra. This is the correct target for the
full non-scalar genus tower.
The intrinsic resolution is Theorem~\ref{thm:mc2-full-resolution}.
\end{definition}

\begin{example}[Heisenberg specialization]
For the Heisenberg algebra $\mathcal{H}_k$ at level $k$,
Theorem~\ref{thm:prism-higher-genus} identifies the genus-$g$ bar
complex $\bar{B}^{(g)}(\mathcal{H}_k)$ with the Feynman transform of
$\mathcal{H}_k$ over the modular operad
$\{\overline{\mathcal{M}}_{g,n}\}$.  Because $\mathcal{H}_k$ has a
single generator with quadratic OPE, each edge contributes the scalar
propagator~$k$.  The diagonal scalar obstruction is
\[
\mathrm{obs}_g(\mathcal H_k)=k\cdot\lambda_g,
\]
and integration against the Faber--Pandharipande class gives the
coefficient computed in~\S\ref{sec:frame-genus-tower}.
\end{example}

\begin{corollary}[The prism principle; \ClaimStatusProvedHere]\label{cor:prism-principle}
The geometric bar complex acts as a mathematical prism decomposing a
chiral algebra $\mathcal{A}$ into its \emph{operadic spectrum}:
\begin{enumerate}[label=\textup{(\roman*)}]
\item \emph{Genus-$0$ spectrum:} The bar complex extracts the
 $\mathrm{Grav}^c$-coalgebra structure on $\bar{B}^{\mathrm{ch}}(
 \mathcal{A})$; the ``colors'' are indexed by the gravity cooperad,
 with OPE structure constants appearing as residues at boundary
 divisors of $\overline{\mathcal{M}}_{0,n}$.
\item \emph{Higher-genus harmonics:} At genus $g \geq 1$, the
 Feynman transform adds ``higher harmonics'' controlled by the
 cohomology of $\overline{\mathcal{M}}_{g,n}$: central extensions
 at genus~$1$, Siegel modular corrections at genus~$2$, and
 so forth.
\item \emph{Reconstruction:} The cobar construction
 $\Omega^{\mathrm{ch}}$ acts as the inverse prism, reconstructing
 $\mathcal{A}$ from its operadic spectrum. For Koszul chiral
 algebras, this reconstruction is exact:
 $\Omega^{\mathrm{ch}}(\bar{B}^{\mathrm{ch}}(\mathcal{A}))
 \xrightarrow{\sim} \mathcal{A}$.
\end{enumerate}
\end{corollary}

\begin{proof}
Part (i) is the content of Theorem~\ref{thm:prism-operadic}(ii):
the bar complex is the operadic bar for the Koszul operad
$\mathrm{HyCom}$, so it produces a $\mathrm{Grav}^c$-coalgebra.
The residue map $d_{\mathrm{res}}$ extracts OPE coefficients, which
are the components of the twisting cochain
$\kappa\colon \mathrm{Grav}^c \to \mathrm{HyCom}$.

Part (ii) follows from Theorem~\ref{thm:prism-higher-genus}(ii):
the Feynman transform includes contributions from all genera.

Part (iii) combines
Theorem~\ref{thm:prism-operadic}(iv)
with Corollary~\ref{cor:bar-cobar-inverse}.
\end{proof}

\begin{example}[Heisenberg specialization]
For $\mathcal{H}_k$, Corollary~\ref{cor:prism-principle} specializes as follows. The genus-$0$ spectrum is the $\mathrm{Grav}^c$-coalgebra encoding the single OPE pole $J(z)J(w) \sim k/(z-w)^2$. The higher-genus harmonics produce the free energy $F_g(\mathcal{H}_k) = k \cdot \frac{2^{2g-1}-1}{2^{2g-1}} \cdot \frac{|B_{2g}|}{(2g)!}$ (where $B_{2g}$ are Bernoulli numbers), recovering the genus expansion of \S\ref{sec:frame-genus-tower}. Reconstruction via the cobar is exact: $\Omega^{\mathrm{ch}}(\bar{B}^{\mathrm{ch}}(\mathcal{H}_k)) \xrightarrow{\sim} \mathcal{H}_k$, since $\mathcal{H}_k$ is Koszul (\S\ref{sec:frame-complementarity}).
\end{example}

\begin{corollary}[Feynman transform involution {\cite{GeK98}}; \ClaimStatusProvedElsewhere]%
\label{cor:feynman-transform-involution}
\index{Feynman transform!involution}
Let $\mathcal{M}od$ be a modular operad and $\mathcal{A}$ an algebra
over $\mathcal{M}od$. Then there is a natural weak equivalence
\begin{equation}\label{eq:ft-involution}
\mathrm{FT} \circ \mathrm{FT}(\mathcal{A})
\;\xrightarrow{\;\sim\;}
\mathcal{A}.
\end{equation}
In particular, the Feynman transform is a homotopy involution on
the category of modular operad algebras.
\end{corollary}

\begin{proof}
This is \cite[Theorem~5.4]{GeK98}. Getzler--Kapranov construct
a natural chain map
$\mathrm{FT}(\mathrm{FT}(\mathcal{A})) \to \mathcal{A}$
and prove it is a quasi-isomorphism by a spectral sequence argument
on the genus filtration: at genus~$0$, the double Feynman transform
reduces to the operadic bar-cobar resolution
$\Omega(\barB(\mathcal{P})) \xrightarrow{\sim} \mathcal{P}$
(which is a quasi-isomorphism for any operad), and the higher-genus
contributions are acyclic by the contractibility of the modular
envelope of the bar-cobar resolution.
\end{proof}

\begin{remark}[Relation to bar-cobar inversion]\label{rem:ft-vs-bar-cobar}
Corollary~\ref{cor:feynman-transform-involution} is the
\emph{all-genera} generalization of our Main Theorem~B
(Theorem~\ref{thm:higher-genus-inversion}): at genus~$0$,
$\mathrm{FT} \circ \mathrm{FT} \simeq \mathrm{Id}$ specializes
to $\Omega \circ \barB \simeq \mathrm{Id}$. The full modular
operad statement applies to the completed all-genera Feynman
transform.  A fixed genus piece
$\bar{B}^{(g)}(\mathcal{A})$ is a graded summand of that completed
object; it need not reconstruct~$\mathcal{A}$ by itself.  The
faithful object is the completed genus tower together with its
Feynman-transform differential and clutching maps.
\end{remark}

\subsection{The modular convolution algebra}
\label{subsec:modular-convolution}
\index{modular convolution dg Lie algebra|textbf}

The Feynman transform $\mathrm{FT}_{\cM\!od}(\cA)$ carries a
modular cooperad coalgebra structure. Its coderivation complex is
the deformation-theoretic dual: a dg~Lie algebra whose MC elements
encode the algebraic structures that the Feynman transform
classifies. This dg~Lie algebra is a \emph{strict model} of the
underlying modular $L_\infty$ convolution algebra
(Definition~\ref{def:modular-convolution-dg-lie}); the MC moduli
coincide, but the full $L_\infty$ structure is needed for transfer,
formality, and gauge equivalence. We define the strict model and
show that its genus-$0$ shadow is the tree-level chiral
convolution.

\begin{definition}[Modular convolution dg~Lie algebra]
\label{def:modular-convolution-vol1}
\index{modular convolution dg Lie algebra!definition}
Let $\cA$ be a cyclic chiral algebra on a smooth projective
curve~$X$ with non-degenerate invariant form. The \emph{modular
convolution dg~Lie algebra} is
\begin{equation}\label{eq:modular-convolution-vol1}
\gAmod
\;:=\;
\operatorname{CoDer}(\mathrm{FT}_{\cM\!od}(\cA))[1]
\;\cong\;
\prod_{\substack{g,n \\ 2g-2+n > 0}}
\Hom_{\Sigma_n}\!\bigl(
C_*(\overline{\cM}_{g,n}),\,
\End_\cA(n)
\bigr),
\end{equation}
completed with respect to the genus grading.
The full dg~Lie structure (differential and bracket) is developed
in Definition~\ref{def:modular-convolution-dg-lie}.
\end{definition}

\begin{theorem}[dg~Lie structure; \ClaimStatusProvedHere]
\label{thm:modular-convolution-structure}
\index{modular convolution dg Lie algebra!structure}
The dg~Lie structure on $\gAmod$ has differential
$D = d_{\mathrm{int}} + [\tau,-] + d_{\mathrm{sew}}
+ d_{\mathrm{pf}} + \hbar_{\mathrm{g}}\Delta$
\textup{(}internal, kernel twist, separating edge, planted-forest
correction, non-separating clutching\textup{)} and Lie bracket
$[\alpha,\beta] = \sum_{\text{sep.\ clutchings}}
\alpha \circ_i \beta
- (-1)^{|\alpha||\beta|} \beta \circ_j \alpha$.
$D^2 = 0$ follows from $\partial^2 = 0$ on
$\overline{\cM}_{g,n}$; Jacobi follows from associativity
of the modular operad composition.
\end{theorem}

\begin{proof}
This is the modular operad convolution of
Getzler--Kapranov~\cite[\S4--5]{GeK98}, applied to
$\cM\!od = \{\overline{\cM}_{g,n}\}$ and the cyclic
algebra~$\cA$. The differential~$D$ transports the boundary
operator~$\partial$ on $C_*(\overline{\cM}_{g,n})$ through the
Hom functor, so $\partial^2 = 0$ gives $D^2 = 0$. The Jacobi
identity follows from the associativity of iterated clutchings.
\end{proof}

\begin{definition}[Chiral convolution dg~Lie algebra]
\label{def:chiral-convolution-vol1}
\index{chiral convolution dg Lie algebra|textbf}
The \emph{chiral convolution dg~Lie algebra} is the tree-level
coderivation complex:
\[
\mathfrak{g}^{\mathrm{ch}}_\cA
\;:=\;
\operatorname{CoDer}(\barB^{\mathrm{ch}}(\cA))[1],
\]
where $\barB^{\mathrm{ch}}(\cA) = \bar{B}^{(0)}(\cA)$ is the
genus-$0$ bar coalgebra.
\end{definition}

\begin{theorem}[Genus completion; \ClaimStatusProvedHere]
\label{thm:vol1-genus-completion}
\index{genus completion}
The genus filtration on $\mathrm{FT}_{\cM\!od}(\cA)$ induces a
complete descending filtration on~$\gAmod$ satisfying:
\begin{enumerate}[label=\textup{(\roman*)}]
\item $\gr_F^g\, \mathrm{FT}_{\cM\!od}(\cA) \cong
 \bar{B}^{(g)}(\cA)$: the genus-$g$ bar complex is the
 genus-$g$ graded piece. At genus~$0$, only trees survive:
 $\bar{B}^{(0)} = \barB_{\mathrm{op}}$
 \textup{(}Corollary~\textup{\ref{cor:feynman-genus0-reduction}}\textup{)}.

\item $\gr_F^0\, \gAmod \cong \mathfrak{g}^{\mathrm{ch}}_\cA$:
 the chiral convolution is the genus-$0$ shadow.

\item $\gAmod = \varprojlim_g\, \gAmod / F^g$:
 the modular convolution is the genus completion of the
 chiral convolution.
\end{enumerate}
The non-separating clutching $\Delta \colon F^g \to F^{g+1}$
raises genus by one and is the sole primitive that creates
genuine modular obstructions.
\end{theorem}

\begin{proof}
The differential decomposes as $D = D_0 + D_1$ where
$D_0$ preserves genus and $D_1 = \hbar_{\mathrm{g}}\Delta$ raises it by one.
The associated graded at genus~$g$ picks out stable graphs of
total genus exactly~$g$. A coderivation of the Feynman transform
restricts, via $F^0/F^1$, to a coderivation of the genus-$0$ bar
coalgebra; the kernel is $F^1\operatorname{CoDer}$, so $\gr^0_F \gAmod =
\mathfrak{g}^{\mathrm{ch}}_\cA$. Completeness follows from
finiteness of stable graphs at each genus.
\end{proof}

The Lie bracket on $\gAmod$ is a correspondence product in the
sense of Chriss--Ginzburg: for $\alpha$ at $(g_1,n_1)$ and $\beta$
at $(g_2,n_2)$, the bracket $[\alpha,\beta]$ composes through the
separating clutching correspondence
$\overline{\cM}_{g_1,n_1+1} \times
\overline{\cM}_{g_2,n_2+1} \xleftarrow{p} Z \xrightarrow{\xi}
\overline{\cM}_{g,n}$
via $[\alpha,\beta] = \xi_* \circ p^*(\alpha \boxtimes \beta)$.
As the Steinberg self-intersection presents the Hecke algebra, the
modular clutching correspondence presents the modular convolution.
The Jacobi identity is the associativity of the correspondence.

\begin{proposition}[The algebra structure as MC element;
\ClaimStatusProvedHere]
\label{prop:vol1-structure-as-MC}
\index{Chriss--Ginzburg structure principle|textbf}
The assignment $\cA \mapsto \Theta_\cA \in \MC(\gAmod)$
is a functor from cyclic chiral algebras to MC moduli of
the modular convolution algebra.
The algebra structure on $\cA$ \emph{is} $\Theta_\cA$: the bar
complex is the twisting of $\gAmod$ by $\Theta_\cA$; the genus
tower is the filtration of $\Theta_\cA$ by graph loop order;
the modular invariants $\kappa$, $\mathfrak{C}$, $\mathfrak{Q}$
are projections of this single element. The genus-$0$ projection
$\Theta_\cA \mapsto \pi$ recovers the universal twisting morphism
of Proposition~\textup{\ref{prop:universal-twisting-adjunction}};
the positive-genus corrections are the modular deformation data.
\end{proposition}

\begin{proof}
The MC equation $D\Theta + \tfrac{1}{2}[\Theta,\Theta] = 0$ is
the genus-expanded nilpotence $D_\cA^2 = 0$
(Theorem~\ref{thm:prism-higher-genus}(iii)). Functoriality holds
because the transferred operations $\ell_n^{\mathrm{tr}}$ and the
propagator $P_\cA$ are natural. The genus filtration of
$\Theta_\cA$ is the graph-loop-order filtration; projection to
degree~$r$ gives the degree-$r$ shadow.
\end{proof}

\begin{remark}[Factorisation homology]
\label{rem:factoris-homology-vol1}
\index{factorisation homology!modular convolution}
In the Beilinson--Drinfeld framework, $\mathfrak{g}^{\mathrm{ch}}_\cA$
is $\int_{\Ran(X)} \End_\cA$ and $\gAmod$ is
$\int_{\cM\!od} \End_\cA$. The MC element $\Theta_\cA$ is a
section of the factorisation algebra over the moduli tower;
its factorisation along clutching maps is the factorisation property.
The ribbon-graph variant ${\gAmod}^{E_1} = \operatorname{CoDer}(F\!\Ass(\cA))[1]$
is the ordered convolution algebra, and its genus expansion is the
algebraic 't~Hooft $1/N$ expansion: planar ribbon graphs dominate
at large~$N$.
\end{remark}

\subsection{Logarithmic forms}

\begin{proposition}[Forced by conformal invariance {\cite{BPZ84}}; \ClaimStatusProvedElsewhere]\label{prop:log-forms-conformal-invariance}
The logarithmic form
$\eta_{ij}=d\log(z_i-z_j)$ is determined by the local conformal
constraints at a collision divisor:
\begin{enumerate}
\item \emph{Translation and scale invariance.} Under
 $z\mapsto az+b$, $\eta_{ij}$ is invariant.
\item \emph{Monodromy control.} A simple logarithmic singularity is
 the pole order compatible with the residue calculus on the boundary
 divisor.
\item \emph{Residue theorem.} The residue of $\eta_{ij}$ along
 $z_i=z_j$ is the unit collision class.
\end{enumerate}
\end{proposition}

\begin{convention}[Signs from trees]
For the bar differential on decorated trees, we use the following sign convention:
\begin{enumerate}
\item Label edges by depth-first traversal starting from the root
\item For contracting edge $e$ connecting vertices with operations $p_1, p_2$ of degrees $|p_1|, |p_2|$:
\item The sign is $(-1)^{\epsilon(e)}$ where:
\[\epsilon(e) = \sum_{e' < e} |p_{s(e')}| + |p_1| + 1\]
where $s(e')$ is the source vertex of edge $e'$ and the sum is over edges preceding $e$ in the ordering.
\item The extra $+1$ comes from the suspension in the bar construction.
\end{enumerate}

% Add missing verification
To verify $d^2 = 0$ for this sign convention, consider a tree with three vertices and two edges $e_1, e_2$. The two ways to contract both edges give:
\begin{itemize}
\item Contract $e_1$ then $e_2$: sign is $(-1)^{\epsilon(e_1)} \cdot (-1)^{\epsilon'(e_2)}$
\item Contract $e_2$ then $e_1$: sign is $(-1)^{\epsilon(e_2)} \cdot (-1)^{\epsilon'(e_1)}$
\end{itemize}
where $\epsilon'$ accounts for the change in edge labeling after the first contraction. A detailed calculation shows these contributions cancel:
\[(-1)^{\epsilon(e_1) + \epsilon'(e_2)} + (-1)^{\epsilon(e_2) + \epsilon'(e_1)} = 0\]
This generalizes to all trees by induction on the number of edges.

This ensures $d^2 = 0$ by a careful analysis of double contractions.
\end{convention}

\begin{lemma}[Sign consistency for bar differential {\cite{LV12}}; \ClaimStatusProvedElsewhere]\label{lem:sign-consistency-bar}
The sign convention above ensures that for any pair of edges $e_1, e_2$ in a tree, the signs arising from contracting $e_1$ then $e_2$ versus contracting $e_2$ then $e_1$ differ by exactly $(-1)$, ensuring $d^2 = 0$.
\end{lemma}

\begin{proof}
This is a standard result in operad theory; see Loday--Vallette \cite{LV12}, Proposition~6.3.2. The key mechanism: label the edges $1, \ldots, E$. Contracting edge $i$ produces a tree with $E-1$ edges, relabeled $1, \ldots, E-1$. If $i < j$, then contracting $i$ first shifts edge $j$ to position $j-1$. The sign for contracting edge at position $k$ in a tree $T$ is $\epsilon(k,T) = k - 1 + \sum_{\ell < k} |v_\ell|$ where $|v_\ell|$ is the degree of the operation at the vertex below edge $\ell$. One checks:
\[\epsilon(i, T) + \epsilon(j-1, T/e_i) = \epsilon(j, T) + \epsilon(i, T/e_j) + 1\]
so that contracting in opposite order introduces a sign of $(-1)^1 = -1$, giving $d^2 = 0$.
\end{proof}

The dual cobar construction is as follows.

\begin{definition}[Cobar construction]
Dually, for a coaugmented cooperad $C$ with coaugmentation $\eta : \mathbb{I} \to C$, the cobar construction $\Omega(C)$ is the free operad on the desuspension $s^{-1}\bar{C}$ (where $\bar{C} = \text{coker}(\eta)$) with differential induced by the cooperad comultiplication.
\end{definition}

\begin{theorem}[Bar-cobar adjunction {\cite{LV12}}; \ClaimStatusProvedElsewhere]\label{thm:bar-cobar-adjunction-operadic}
There is an adjunction:
\[
\barB : \text{Operads} \rightleftarrows \text{Cooperads}^{\text{op}} : \Omega
\]
If $P$ is Koszul (Definition~\ref{def:koszul-operad}), then the unit and counit are quasi-isomorphisms, establishing an equivalence of homotopy categories.
\end{theorem}

\subsection{Partition complexes and the commutative operad}

For the commutative operad $\Com$, the bar construction admits an explicit combinatorial model via partition lattices:

\begin{definition}[Partition lattice]
The partition lattice $\Pi_n$ is the poset of all partitions of $\{1, 2, \ldots, n\}$, ordered by refinement: $\pi \leq \sigma$ if every block of $\pi$ is contained in some block of $\sigma$. The proper part $\barPi_n = \Pi_n \setminus \{\hat{0}, \hat{1}\}$ excludes the minimum (discrete partition) and maximum (trivial partition).
\end{definition}

\begin{theorem}[Partition complex structure; \ClaimStatusProvedHere]\label{thm:partition}
The bar complex $\barB(\Com)(n)$ is quasi-isomorphic to the reduced chain complex $\tilde{C}_*(\barPi_n)$ of the proper part of the partition lattice $\Pi_n$. More precisely:
\[
\barB(\Com)(n) \simeq s^{n-1}\tilde{C}_{*}(\barPi_n) \otimes \sgn_n
\]
where $\sgn_n$ is the sign representation of $S_n$.
\end{theorem}

\begin{proof}
Elements of $\Com^{\circ k}(n)$ (the $k$-fold composition) correspond to ways of iteratively partitioning $n$ elements through $k$ levels. The simplicial structure is:
\begin{itemize}
\item Face maps compose adjacent levels of partitioning (coarsening)
\item Degeneracy maps repeat a level (refinement followed by immediate coarsening)
\end{itemize}

After normalization (removing degeneracies), we obtain chains on $\barPi_n$. The dimension shift and sign representation arise from the suspension in the bar construction and the need for $S_n$-equivariance.

The poset $\barPi_n$ has the homology of a wedge of $(n-1)!$ spheres of dimension $n-3$, with the $S_n$-action on the top homology given by the Lie representation tensored with the sign. This follows from the classical results of Björner--Wachs \cite{BW83} and Stanley \cite{Sta97}, who computed:
\[
\tilde{H}_{n-3}(\barPi_n) \cong \Lie(n) \otimes \sgn_n \text{ as } S_n\text{-representations}
\]
and $\tilde{H}_k(\barPi_n) = 0$ for $k \neq n-3$.
\end{proof}
\begin{remark}[Simplicial model]
The simplicial bar for $\Com$ literally consists of chains of refinements $\pi_0 \leq \pi_1 \leq \cdots \leq \pi_k$ in $\Pi_n$. This is the nerve of the poset $\Pi_n$, and the identification with the cooperad structure follows from taking normalized chains.
\end{remark}

\subsection{Holographic interpretation}

\begin{conjecture}[Holographic Hochschild comparison; \ClaimStatusConjectured]\label{conj:holographic-koszul}
Let $\cA_{\mathrm{bdy}}$ be the boundary chiral algebra of a
$3$d holomorphic-topological theory.  The algebraic bulk attached to
the boundary is
\[
Z^{\mathrm{der}}_{\mathrm{ch}}(\cA_{\mathrm{bdy}})
=C^\bullet_{\mathrm{ch}}(\cA_{\mathrm{bdy}},\cA_{\mathrm{bdy}}).
\]
The physical $3$d bulk is the BRST observable algebra
$\Obs_{\mathrm{BRST}}(\mathcal A_{\mathrm{bulk}}^{3d})$ of the
AdS$_3$-side theory.  The holographic conjecture is the existence,
after the higher-spin target has been constructed and compared with
the BRST complex, of an $E_2$-algebra equivalence
\[
Z^{\mathrm{der}}_{\mathrm{ch}}(\cA_{\mathrm{bdy}})
\;\simeq\;
\Obs_{\mathrm{BRST}}(\mathcal A_{\mathrm{bulk}}^{3d})
\]
compatible with the boundary action.  Boundary observables and
$3$d HT amplitudes are then related by the square below.

\begin{center}
\begin{tikzcd}
\text{Boundary chiral algebra } \cA_{\mathrm{bdy}}
\arrow[r, "Z^{\mathrm{der}}_{\mathrm{ch}}"]
\arrow[d, "\text{correlators}"] &
Z^{\mathrm{der}}_{\mathrm{ch}}(\cA_{\mathrm{bdy}})
\arrow[r, dashed, "\mathrm{BV/BRST}"]
\arrow[d, "\text{bulk action}"] &
\Obs_{\mathrm{BRST}}(\mathcal A_{\mathrm{bulk}}^{3d})
\arrow[d, "\text{amplitudes}"] \\
\text{Boundary observables}
\arrow[r, "\text{open--closed}"] &
\text{Intrinsic bulk operators}
\arrow[r, dashed, "\text{AdS/CFT}"] &
\text{3d bulk amplitudes}
\end{tikzcd}
\end{center}

Specifically:
\begin{enumerate}[label=\textup{(\roman*)}]
\item The boundary-to-intrinsic-bulk map is the centre map
$\cA_{\mathrm{bdy}}\to Z^{\mathrm{der}}_{\mathrm{ch}}(\cA_{\mathrm{bdy}})$,
not the bar construction.
\item The bar coalgebra
$\barBch_X(\cA_{\mathrm{bdy}})$ classifies twisting morphisms and MC
couplings.  Its Verdier dual gives the defect/Koszul-dual lane
$\cA_{\mathrm{bdy}}^!$ after completed cobar.
\item Collision residues model bulk interactions only after passage
to brace operations on chiral Hochschild cochains.
\item Cobar reconstruction recovers the boundary algebra from its bar
coalgebra on the Koszul or completed Positselski surface.  It does not
recover the boundary from the physical $3$d bulk without the extra
holographic comparison.
\item Scalar genus corrections split into the Theorem~D scalar lane
and the cross-channel term
$\delta F_g^{\mathrm{cross}}$ for multi-weight families.  A Witten
diagram computation may match both terms, but the algebraic theorem
does not identify them by itself.
\end{enumerate}

\emph{Example.} For
$\mathcal{A}_{\mathrm{bdy}} = \mathcal{W}_{\infty}[\lambda]$
at $c = N$:
\begin{itemize}
\item Physical $3$d bulk candidate: Vasiliev higher-spin gravity in AdS$_3$.
\item Derived center: should organize the same higher-spin interaction
 data as $\mathrm{hs}[\lambda]\otimes\mathcal{C}^{\bullet}(\mathrm{AdS}_3)$
 after the BV/BRST comparison.
\item Bar complex: provides the coalgebraic coupling input, not the
 closed-sector observable algebra.
\item Cobar complex: recovers the boundary $\mathcal{W}_{\infty}$
 package from the boundary bar coalgebra under the completed
 bar-cobar theorem.
\item The parameter $\lambda$ controls both
 boundary $\cW$-algebra structure constants and
 bulk higher-spin coupling constants.
\end{itemize}
\end{conjecture}

\begin{remark}[Physical evidence and scope]
This conjecture is supported by matching of partition functions,
three-point functions, and conformal blocks between boundary W-algebras
and bulk Vasiliev theory \cite{GG11}. The algebraic theorem supplies
the completed bar-cobar adjunction and its finite-stage coefficient
identities. The holographic assertion adds the identification of
cochain-level chiral closed-sector operators with BRST bulk observables, together
with a comparison between bar-side MC couplings and Witten-diagram
couplings. Those identifications are the AdS$_3$/CFT$_2$ input.
\end{remark}

\begin{remark}[Algebraic core and external comparison]
The algebraic core of Conjecture~\ref{conj:holographic-koszul} is the
completed bar-cobar package. Its finite stages give the coefficient
identities used by the boundary $\mathcal W$-algebra.  A filtered
higher-spin target and the comparison with Witten diagrams and bulk
fields are additional holographic data. Thus the conjecture is not
used as input to the bar-cobar theorem; it is an application of the
theorem after the Gaberdiel--Gopakumar and AdS$_3$/CFT$_2$
identifications are supplied.
\end{remark}

\section{K3 Mukai Interface}
\label{sec:pdq-supplement}

\begin{remark}[Mukai pairing $(4,20)$ as quantum Poincar\'e duality]
\label{rem:pdq-1}
The Mukai pairing on
$\Lambda_{\mathrm{Muk}} = H^\star(K3,\mathbb Z) \cong II_{4,20}$
is the Calabi--Yau-side form whose chiral image is the
quantum-Poincar\'e-duality pairing on $\mathbf H_{\Delta_5}$.  On the
chiral side the bulk pairing is the Hochschild/Mukai pairing on
$\ChirHoch^\bullet(\mathbf H_{\Delta_5})$; the Koszul-dual line is the
Verdier defect lane $\mathbf H_{\Delta_5}^!$.
The $(4,20)$ signature is triply visible:
\begin{enumerate}
 \item Mukai-doubling $K^{\kappa_{\mathrm{ch}}} = 2c_+ = 8$;
 \item Shadow-tower depth at the class-$\mathcal M$ Virasoro face
 $c_{2d} = -214 = -2\cdot 107$ with $107 = 5\cdot 24 - 13$
 (Shapere--Tachikawa 2008 \S 3, $c_{4d} = (5n-13)/6$ at $n=24$);
 \item Andrianov factorisation shifts
 $L(s,\Phi_{10}) = L(s,\Delta\otimes E_6)\,\zeta(s-9)\,\zeta(s-8)$,
 with $9 = c_-/2-1$ and $8 = K^{\kappa_{\mathrm{ch}}}$.
\end{enumerate}
The $\kappa^!$-complementarity between $\kappa_{\mathrm{ch}}, \kappa_{\mathrm{cat}}, \kappa_{\mathrm{BKM}}, \kappa_{\mathrm{fiber}}$
is the K3-chiral-bialgebra refinement of derived-centre complementarity
\cite{Mukai1987,Bruinier2002,Borcherds1998}.
Vol.~III Chapter~\ref{V3-ch:k3-chiral-bialgebra-platonic} records
the K3 chiral-bialgebra chart.  The chain-level realisation of the
Mukai form as a non-degenerate
pairing on $\ChirHoch^\bullet(\mathbf H_{\Delta_5})$ is
Theorem~\ref{thm:chi3-mukai-hochschild-pairing}
(\texttt{hochschild\_cohomology.tex}~\S~\ref{sec:chirhoch-mukai-nc-hodge}):
on the Koszul locus $\Uadm\subset\overline{\mathcal A_2}$ the
chiral non-commutative Hodge-to-de-Rham spectral sequence
(Kaledin~$2017$~\emph{Duke}~$166$~Thm.~$5.4$ transported via
Francis--Gaitsgory smoothness) degenerates at~$E_1$, inducing a
CY-$3$ Serre autoequivalence $\sSer$ whose pairing evaluates to
$\langle[\chi_3],[\chi_3^\vee]\rangle^{\mathrm{ch}}_{\Muk}
=2\,\mathrm{Vol}(E)\cdot(2\pi\mathrm i)^3$, coincident with the
Siegel--Eisenstein period
$\chi(\mathcal O_{\mathrm{K3}})\cdot
\mathrm{Res}_{s=1/2}E_5^{(2)}|_{K(1)}$ of
Theorem~\ref{thm:chi3-siegel-eisenstein-period}.
\end{remark}

\begin{remark}[Shadow-tower depth at the class-$\mathcal M$ Virasoro face]
\label{rem:pdq-2}
On $\mathcal T[A_1,\Sigma_{0,24}]$: Coulomb rank $21$,
$c_{4d} = 107/6 = (5n-13)/6|_{n=24}$ (Shapere--Tachikawa 2008
\S 3 applied to Chacaltana--Distler 2010 eq.~(5.14) count
$(n_v, n_h) = (63, 88)$), $c_{2d} = -12 c_{4d} = -214$
(Beem--Rastelli, BLLPRvR 2013 eq.~(3.14)). The Virasoro face of the
shadow tower is $c_{2d} = -2\cdot 107$ with $107$ prime.  The
class-$\mathcal S$ computation gives the first-principles structure
$5n-13$ at $n=24$; the number is not obtained by decomposing the
Mukai signature. The quantum-Poincar\'e-duality reads
this as a chiral/vertex-algebra Poincar\'e dual of the Mukai
pairing: class-$\mathcal M$ Virasoro on the chiral side dualises
the Mukai $(4,20)$ on the CY side; the specific central-charge
value within the class-$\mathcal M$ stratum is the
first-principles $-214$.
The class-$\mathcal S$ AGT chiral algebra is recorded in
Vol.~III Chapter~\ref{V3-ch:k3-chiral-bialgebra-platonic}; the
open/closed register used here is
Section~\ref{sec:thqg-open-closed-realization}.
\end{remark}

\begin{remark}[DMVV and Nakajima-BPS character of
$\mathbf{H}_{\Delta_{5}}$]
\label{rem:pdq-witten-dmvv-nakajima}
\index{DMVV second-quantised genus}%
\index{Nakajima!Heisenberg on Hilbert scheme of K3}%
\index{BPS state count!K3 elliptic genus}%
\index{Witten!BPS / second-quantisation}%
The second-quantised symmetric-product formula of
Dijkgraaf--Moore--Verlinde--Verlinde 1997
(\texttt{hep-th/9608096}, Theorem~3),
\[
 \sum_{n \geq 0}\,p^{n}\,Z_{\mathrm{Sym}^{n}(K3)}(\tau, z)
 \;=\;
 \prod_{m > 0,\, n \geq 0,\, \ell \in \mathbb{Z}}
 \bigl(1 - p^{m}\,q^{n}\,y^{\ell}\bigr)^{-c(mn, \ell)},
\]
with $c(m n, \ell)$ the Fourier coefficients of the K3 elliptic genus
$\phi_{\mathrm{K3}}(\tau, z) = 2\phi_{0,1}(\tau, z)$ (Kawai--Yamada--Yang
1994; Eguchi--Ooguri--Tachikawa 2011), coincides with the Borcherds
multiplicative lift $1/\Phi_{10}^{\mathrm{un}}(Z)$ in variables
$p = e^{2\pi i \sigma}$, $q = e^{2\pi i \tau}$, $y = e^{2\pi i z}$
(Borcherds 1998 \S10; the master identification traces to
Gritsenko--Nikulin 1998 \texttt{alg-geom/9712007} Theorem~4.2). Vol~III
Theorem~\ref{V3-thm:k3-siegel-character} shows that the genus-$2$ Siegel
character of the K3 chiral bialgebra
$\mathbf{H}_{\Delta_{5}}$
(Vol~III Chapter~\ref{V3-ch:k3-chiral-bialgebra-platonic}
Definition~\ref{V3-def:k3-chiral-bialgebra}) is $1/\Phi_{10}^{\mathrm{un}} =
1/\Delta_{5}^{\,2}$, so the DMVV identity realises the quantum
Poincar\'e-duality pairing of this section as the generating function
of symmetric-product-K3 Hilbert spaces. The Verdier defect lane
\[
\Omegach_X\mathcal{D}_X(\mathbf{H}_{\Delta_{5}})
\simeq
\mathbf{H}_{\Delta_{5}}^{\,!}
\]
bosonises this BPS generating function:
on $\mathrm{Sym}^{n}(K3)$ the Hilbert space is
$V_{\Lambda_{\mathrm{Muk}}}^{\otimes n}/S_{n}$
(the abelian-Heisenberg Fock from the orbifold symmetric product), while
on $\mathrm{Hilb}^{n}(K3)$ the Nakajima enhancement (Nakajima 1997
\emph{Lectures on Hilbert Schemes of Points on Surfaces} AMS ULS~18
Theorem~8.13; independently Grojnowski 1996
\texttt{alg-geom/9506020}) promotes the abelian Fock to a Heisenberg
algebra whose raising and lowering operators generate the non-abelian
BKM $\mathfrak{g}_{\Delta_{5}}$ by the Frenkel--Kac mechanism. The
Mukai--Fourier--Mukai kernel
$\Delta^{\vee}_{\mathrm{Muk}} \in D^{b}(K3 \times K3)$ is the
CY-side kernel whose chiral transform gives the Hochschild/Mukai
pairing on $Z_{\mathrm{der}}^{\mathrm{ch}}(\mathbf{H}_{\Delta_{5}})$.
It is not the bar-cobar counit and not the bar coalgebra; its
comparison with the Verdier defect lane uses
Theorem~\ref{thm:defect-koszul} together with the Hochschild pairing
theorem.  The resulting intrinsic-bulk pairing is the algebraic side
of the holographic comparison with the BPS microstate count of
Maldacena--Moore--Strominger 1999 \texttt{hep-th/9903163} for the
D1-D5 system on $K3 \times S^{1}$. The ``non-abelian'' attribute of
$\mathbf{H}_{\Delta_{5}}$ (Vol~III Theorem~\ref{V3-thm:k3-abelian-at-lie})
is thus the Nakajima enhancement: abelian at the symmetric-product
Fock level, non-abelian at the Hilbert-scheme vertex-operator closure.
Primary: Dijkgraaf--Moore--Verlinde--Verlinde 1997; Nakajima 1997;
Grojnowski 1996; Maldacena--Moore--Strominger 1999; Borcherds 1998;
Kawai--Yamada--Yang 1994; Dabholkar--Harvey 1989
\texttt{hep-th/8909030} (precursor to the BPS generating function).
Cross-volume bridge: Vol~III
\texttt{chapters/examples/k3\_chiral\_bialgebra\_platonic.tex}
Theorem~\ref{V3-thm:bkm-simple-roots-walls}; Vol~II
\texttt{chapters/connections/modular\_pva\_quantization\_frontier.tex}
Remark~\ref{V2-rem:witten-het-iia-mukai-narain} (Het/IIA duality frame).
\end{remark}
